\chapter{Summary}

This thesis investigated topics of improving FAIRness, accessibility, and reducing barriers. This was conducted in pairs of scientific and technical achievements, across several scales of investigation; from simple prokaryotes in a classroom, through human cell types in a hospital, to applications for humans in the wider world.

In \cref{chapter:annotation} publication presented a working implementation of the GMOD suite of genomics infrastructure for the application of teaching undergraduate bacteriophage genomics courses. Over 50 new or updated tools were used across three workflows, to annotate more than 160+ phages to date over several years of testing in undergraduate environments. Additionally a wrapper for the Circos plotter was presented which significantly improves accessibility and discoverability of Circos and eases its use in reproducible scientific workflow.
In \cref{chapter:clinic} this infrastructure was built upon for the analysis of RNA sequencing data collected from endothelial cells and pericytes exposed to TNF. We discovered that TNF has a significant impact on both endothelial cells and pericytes, and that factors produced by pericytes---when exposed to TNF---have a significant impact on endothelial cells permeability with implications for optimising drug delivery in tumours. In the technical application we explored securing access to sensitive clinical datasets and removing barriers for researchers by providing a more integrated and user-friendly access method.
\Cref{chapter:training} finally looked to dissemination of these important results, how best to ensure sustainability for newly developed tools and resources. We presented a series of improvements to the Galaxy Training Network to improve FAIRness, increase accessibility, and remove barriers to access. A number of improvements were made to directly impact course didactics and improve teachers' experiences such as automatically generated Jupyter notebook for programming lessons and significantly overhauled accessibility via a thorough screenreader testing program. The presented small improvements have big impacts, affecting 45-50k visitors per month. Additionally the process of teaching was made more concrete through the development of Training Infrastructure as a Service (TIaaS) which has provided training infrastructure for 25k students over 5 years.


